% Replacement text for the theory part of Job_Market_paper_Kozyrev-2.pdf.
% Suggested use:
% 1. Replace Section 4.7 with Part I below.
% 2. Replace the current Appendix F/G material on exact bounds and rates with Part II below.

% --------------------------------------------------------------------------
% Part I. Main-text replacement for Section 4.8
% --------------------------------------------------------------------------

\subsection{Bounds on $R_k$}
\label{sec:rk_bounds}

The representation above shows that $R_k^{\mathrm{bag}}$ is the squared
distance between the infinite-bagged RSM weights and the oracle weights.
Although the conditional harmonic means inside $g_i(k)$ do not admit a
general closed form, the excess risk can still be bounded.  Let
\[
  z_k=\frac{1}{k}-\frac{1}{m}.
\]

\begin{proposition}[Exact bracket for infinite-bagged excess risk]
\label{prop:exact_bracket}
For every $1\le k\le m$,
\[
  \frac{
    \left(kE_S[1/A_S]-1/\bar\tau\right)^2
  }{
    \sum_{i=1}^m\sigma_i^2-m/\bar\tau
  }
  \le
  R_k^{\mathrm{bag}}
  \le
  E_S[1/A_S]-\frac{1}{A_m}.
\]
Using the Taylor expansion of $E_S[1/A_S]$, the bracket becomes, to leading
order,
\[
  \frac{V_\delta^2}{m(H-1)\bar\tau}z_k^2
  \lesssim
  R_k^{\mathrm{bag}}
  \lesssim
  \frac{1}{\bar\tau}z_k,
\]
up to the finite-$m$ factor $m^2/(m-1)^2$ in the lower bound.
\end{proposition}

\noindent
The upper bound is the Jensen ceiling and is linear in $z_k$.  The lower
bound is the Cauchy--Schwarz projection floor and is quadratic in $z_k$.
Thus the two boundary criteria correspond to the two endpoint rates:
the linear ceiling gives the square-root rule, whereas the quadratic floor
gives the cube-root rule.  A level bracket alone, however, does not pin down
the minimizer of the true loss: for that we need a regularity condition on
the one-step decline of $R_k^{\mathrm{bag}}$.

\begin{assumption}[Marginal regularity]
\label{ass:marginal_regular}
Let
\[
  \Delta R_k=R_k^{\mathrm{bag}}-R_{k+1}^{\mathrm{bag}}.
\]
Consider a triangular array indexed by $T$, with pool size $m=m_T$ and
excess-risk sequence $R_{k,T}^{\mathrm{bag}}$.  For $1\ll k\ll m_T$,
suppose there exist constants $0<d_L\le d_U<\infty$ that are uniform in
$T$, $m_T$, and $k$, such that
\[
  d_L\left(z_k^2-z_{k+1}^2\right)
  \le
  \Delta R_k
  \le
  d_U\left(z_k-z_{k+1}\right).
\]
Equivalently, for $k\ll m$,
\[
  \Delta R_k\gtrsim k^{-3},
  \qquad
  \Delta R_k\lesssim k^{-2}.
\]
\end{assumption}

\begin{proposition}[Rate window for the optimal subset size]
\label{prop:rate_window}
Let
\[
  L_k^{\mathrm{bag}}
  =
  \left(A+R_k^{\mathrm{bag}}\right)\frac{T-1}{T-k},
  \qquad
  A=\sigma_\epsilon^2+Q^{\mathrm{opt}}>0,
\]
and let
\[
  k_T^*\in\arg\min_{1\le k\le m}L_k^{\mathrm{bag}}.
\]
Suppose Assumption~\ref{ass:marginal_regular} holds uniformly along the
triangular array, $R_{k,T}^{\mathrm{bag}}$ is uniformly bounded, and the
pool size is large enough that the relevant range satisfies
$T^{1/2}=o(m_T)$.  Then
\[
  k_T^*=\Omega(T^{1/3}),
  \qquad
  k_T^*=O(T^{1/2}).
\]
Thus the optimal subset size lies between the cube-root and square-root
rates:
\[
  T^{1/3}\lesssim k_T^*\lesssim T^{1/2}.
\]
If $m$ is fixed or grows more slowly than $T^{1/2}$, the upper rate is
truncated by the constraint $k\le m$.
\end{proposition}

\noindent
The uniform marginal condition is substantive.  It is not implied by the
level bracket in Proposition~\ref{prop:exact_bracket}; in particular, the
coefficient in the displayed quadratic lower envelope contains $1/m$ and
may vanish when the forecast pool grows.  The proposition should therefore
be read as a conditional rate result, not as an unconditional consequence
of the common-loading model.  With a fixed pool, no nontrivial
$T\to\infty$ rate is possible because the optimizer is eventually capped
at $m$.

The proofs of Propositions~\ref{prop:exact_bracket} and
\ref{prop:rate_window} are given in Appendix~\ref{app:rk_bounds_rates}.

% --------------------------------------------------------------------------
% Part II. Appendix replacement for current Appendix F/G material
% --------------------------------------------------------------------------

\section{Exact Bounds on $R_k$ and the Optimal Subset-Size Window}
\label{app:rk_bounds_rates}

\subsection{Proof of Proposition~\ref{prop:exact_bracket}}

The upper bound follows from Jensen's inequality.  Since
$v\mapsto v'\Sigma_e v$ is convex,
\[
  Q_k^{\mathrm{bag},\infty}
  =
  \left(E_S v_S\right)'\Sigma_e\left(E_S v_S\right)
  \le
  E_S[v_S'\Sigma_e v_S]
  =
  Q_k^{\mathrm{sub}}.
\]
Subtracting $Q^{\mathrm{opt}}$ and using
\[
  Q_k^{\mathrm{sub}}=q+E_S[1/A_S],
  \qquad
  Q^{\mathrm{opt}}=q+\frac{1}{A_m},
\]
gives
\[
  R_k^{\mathrm{bag}}
  \le
  E_S[1/A_S]-\frac{1}{A_m}.
\]

For the lower bound, define
\[
  h_i(k)=g_i(k)-\frac{1}{\bar\tau}.
\]
The score identities imply
\[
  \sum_{i=1}^m\tau_i g_i(k)=m,
  \qquad
  \sum_{i=1}^m g_i(k)=kmE_S[1/A_S].
\]
Therefore
\[
  \sum_{i=1}^m\tau_i h_i(k)=0
\]
and
\[
  \sum_{i=1}^m h_i(k)
  =
  m\left(kE_S[1/A_S]-\frac{1}{\bar\tau}\right)
  \equiv d_k.
\]
The true vector $h(k)$ is feasible for the constrained minimization problem
\[
  \min_h \sum_{i=1}^m\tau_i h_i^2
  \quad\text{s.t.}\quad
  \sum_i\tau_i h_i=0,\qquad
  \sum_i h_i=d_k.
\]
The Lagrangian first-order condition implies
\[
  h_i=a+\frac{b}{\tau_i}.
\]
Using $\sum_i\tau_i h_i=0$ gives $a=-b/\bar\tau$, hence
\[
  h_i=b\left(\frac{1}{\tau_i}-\frac{1}{\bar\tau}\right).
\]
Using $\sum_i h_i=d_k$ gives
\[
  b=
  \frac{d_k}{
    \sum_i1/\tau_i-m/\bar\tau
  }.
\]
Thus the minimum weighted norm is
\[
  \sum_i\tau_i h_i^2
  =
  \frac{d_k^2}{
    \sum_i1/\tau_i-m/\bar\tau
  }.
\]
Since
\[
  R_k^{\mathrm{bag}}
  =
  \frac{1}{m^2}\sum_i\tau_i h_i(k)^2
\]
and $d_k=m(kE_S[1/A_S]-1/\bar\tau)$, we obtain
\[
  R_k^{\mathrm{bag}}
  \ge
  \frac{
    \left(kE_S[1/A_S]-1/\bar\tau\right)^2
  }{
    \sum_i1/\tau_i-m/\bar\tau
  }.
\]
Finally, $1/\tau_i=\sigma_i^2$, giving the stated lower bound.

\subsection{Taylor Expansion in $z_k$}
\label{app:zk_expansion}

The Taylor expansion of $E_S[1/A_S]$ around $k\bar\tau$ is
\[
  E_S[1/A_S]
  =
  \frac{1}{k\bar\tau}
  +
  \frac{\operatorname{Var}_S(A_S)}{(k\bar\tau)^3}
  +
  o(z_k).
\]
Under simple random sampling without replacement,
\[
  \operatorname{Var}_S(A_S)
  =
  k\frac{m-k}{m-1}
  \frac{1}{m}\sum_{i=1}^m(\tau_i-\bar\tau)^2
  =
  k^2z_k\frac{m}{m-1}V_\tau,
\]
where
\[
  V_\tau=\frac{1}{m}\sum_{i=1}^m(\tau_i-\bar\tau)^2,
  \qquad
  V_\delta=\frac{V_\tau}{\bar\tau^2}.
\]

For the upper bound,
\[
  E_S[1/A_S]-\frac{1}{A_m}
  =
  \frac{1}{k\bar\tau}-\frac{1}{m\bar\tau}
  +o(z_k)
  =
  \frac{1}{\bar\tau}z_k+o(z_k).
\]
Thus the Jensen ceiling is linear in $z_k$:
\[
  R_k^{\mathrm{bag}}\lesssim \frac{1}{\bar\tau}z_k.
\]

For the lower bound, define
\[
  D_k=kE_S[1/A_S]-\frac{1}{\bar\tau}.
\]
The first-order term cancels, and the Taylor expansion gives
\[
  D_k
  =
  \frac{\operatorname{Var}_S(A_S)}{k^2\bar\tau^3}
  +o(z_k)
  =
  \frac{m}{m-1}\frac{V_\delta}{\bar\tau}z_k
  +o(z_k).
\]
Moreover,
\[
  \sum_i\sigma_i^2-\frac{m}{\bar\tau}
  =
  \frac{m(H-1)}{\bar\tau},
  \qquad
  H=\overline{\sigma^2}\bar\tau.
\]
Therefore the projection floor is quadratic in $z_k$:
\[
  R_k^{\mathrm{bag}}
  \gtrsim
  \frac{V_\delta^2}{m(H-1)\bar\tau}z_k^2,
\]
up to the finite-$m$ factor $m^2/(m-1)^2$.  Hence, to leading order,
\[
  \frac{V_\delta^2}{m(H-1)\bar\tau}z_k^2
  \lesssim
  R_k^{\mathrm{bag}}
  \lesssim
  \frac{1}{\bar\tau}z_k.
\]

\subsection{Boundary Optima and Proof of Proposition~\ref{prop:rate_window}}

We first compute the two boundary optimizers.  For the upper boundary, set
\[
  R_k^U=Bz_k,
  \qquad
  B=\frac{1}{\bar\tau}.
\]
Since $A=\sigma_\epsilon^2+q+B/m$, the loss becomes
\[
  \left(\sigma_\epsilon^2+q+\frac{B}{k}\right)
  \frac{T-1}{T-k}.
\]
The first-order condition gives
\[
  k_U^*
  =
  \frac{
    -B+\sqrt{B^2+(\sigma_\epsilon^2+q)BT}
  }{
    \sigma_\epsilon^2+q
  }
  \sim
  \left(
    \frac{BT}{\sigma_\epsilon^2+q}
  \right)^{1/2}.
\]
Thus the Jensen ceiling gives the square-root rule,
\[
  k_U^*\asymp T^{1/2}.
\]

For the lower boundary, set
\[
  R_k^L=C_Lz_k^2,
  \qquad
  C_L=\frac{V_\delta^2}{m(H-1)\bar\tau}.
\]
The continuous first-order condition is
\[
  -\frac{2C_Lz_k}{k^2}(T-k)+A+C_Lz_k^2=0.
\]
For the interior range $k\ll m$, so that $z_k\approx1/k$, this becomes
\[
  Ak^3+3C_Lk-2C_LT=0.
\]
Hence
\[
  k_L^*
  \sim
  \left(
    \frac{2C_LT}{A}
  \right)^{1/3}
  =
  \left(
    \frac{2V_\delta^2T}
    {m(H-1)\bar\tau A}
  \right)^{1/3}.
\]
Thus the projection floor gives the cube-root rule,
\[
  k_L^*\asymp T^{1/3}.
\]

The level bounds alone do not mechanically imply that the minimizer of the
true $R_k^{\mathrm{bag}}$ lies between these two optima.  For that conclusion
we use Assumption~\ref{ass:marginal_regular}.

The loss decreases from $k$ to $k+1$ if and only if
\[
  L_{k+1}^{\mathrm{bag}}<L_k^{\mathrm{bag}}.
\]
Using
\[
  L_k^{\mathrm{bag}}
  =
  \left(A+R_k^{\mathrm{bag}}\right)\frac{T-1}{T-k},
\]
cross-multiplication gives the equivalent condition
\[
  (T-k)\Delta R_k>A+R_k^{\mathrm{bag}}.
\]

For the lower rate bound, the marginal lower envelope implies
\[
  (T-k)\Delta R_k\gtrsim Tk^{-3}
\]
for $k\ll T$.  Since $A+R_k^{\mathrm{bag}}$ is uniformly bounded above,
the inequality still holds for all $k\le cT^{1/3}$ for some constant
$c>0$.  Hence the loss is decreasing up to the cube-root scale, so the
minimizer cannot occur below it:
\[
  k_T^*=\Omega(T^{1/3}).
\]

For the upper rate bound, the marginal upper envelope implies
\[
  (T-k)\Delta R_k\lesssim Tk^{-2}.
\]
Since $A+R_k^{\mathrm{bag}}\ge A>0$, this quantity is smaller than
$A+R_k^{\mathrm{bag}}$ for all $k\ge CT^{1/2}$, for some sufficiently large
constant $C$.  Hence the loss is increasing beyond the square-root scale,
so
\[
  k_T^*=O(T^{1/2}).
\]
